% Independently reviewed restoration from the July 2026 Encyclopedia.
% Requires the canonical normalized maps, binary contact criterion and modulus
% bounds already established in the main manuscript. No new global shell
% regularity or monotonicity is assumed.
\section{Convex support certificates and rigidity at the outer bound}
\label{sec:restored-support}

Write the normalized similarities as
\[
 f_s(z)=s+L_s z,\qquad s\in\{-1,1\},\qquad L_s=r_sO_s,
\]
where $O_s$ is an orthogonal real-linear map and $0<r_s<1$.
Put $q=\max(r_-,r_+)$, let $K$ be the attractor, and let
$P=\operatorname{conv}K$. Inner products below are the Euclidean inner
product on $\mathbb C\simeq\mathbb R^2$, with $1$ denoting the real unit
vector.

\begin{theorem}[Convex support certificate]
\label{thm:restored-support}
The support function $h(u)=\sup_{x\in P}\langle u,x\rangle$, $u\in S^1$,
is the unique continuous solution of
\begin{equation}\label{eq:restored-support-fixed}
 h(u)=\max_{s\in\{-1,1\}}
 \bigl(s\langle u,1\rangle+r_s h(O_s^{\mathsf T}u)\bigr).
\end{equation}
Define
\begin{equation}\label{eq:restored-support-gap}
 g(u)=r_-h(O_-^{\mathsf T}u)+r_+h(-O_+^{\mathsf T}u)
       -2\langle u,1\rangle.
\end{equation}
Then $f_-(P)\cap f_+(P)\ne\varnothing$ if and only if
$g(u)\ge0$ for every $u\in S^1$. In particular, one direction for which
$g(u)<0$ certifies that $K$ is disconnected.

Let $T$ denote the operator in \eqref{eq:restored-support-fixed}, and set
$h_0=0$, $h_n=T^nh_0$. Then
\begin{equation}\label{eq:restored-support-error}
 0\le h_n\le h_{n+1}\le h,\qquad
 \|h-h_n\|_\infty\le E_n:=\frac{q^n}{1-q}.
\end{equation}
If $g_n$ is obtained by replacing $h$ with $h_n$ in
\eqref{eq:restored-support-gap}, then
\begin{equation}\label{eq:restored-support-gap-error}
 g_n(u)\le g(u)\le g_n(u)+(r_-+r_+)E_n.
\end{equation}
Consequently, $g_n(u)+(r_-+r_+)E_n<0$ is a finite disconnection
certificate.
\end{theorem}

\begin{proof}
Taking the convex hull of the invariance equation gives
$P=\operatorname{conv}(f_-(P)\cup f_+(P))$. The support function of an
affine image is obtained by applying the transpose of its linear part
to the direction and adding the translation term. This proves
\eqref{eq:restored-support-fixed}. The operator $T$ is a contraction on
$C(S^1)$ in the supremum norm, with contraction constant at most $q$,
so its fixed point is unique.

The two convex pieces meet precisely when
$2\in L_-P-L_+P$. Membership in a compact convex set is equivalent to
all its supporting half-space inequalities. The support function of
$L_-P-L_+P$ is
$r_-h(O_-^{\mathsf T}u)+r_+h(-O_+^{\mathsf T}u)$, proving the criterion.
Actual first-piece contact implies convex first-piece contact.

The operator $T$ preserves pointwise order, and
$h_1(u)=|\langle u,1\rangle|\ge0$. Thus its iterates increase to $h$.
Since $\|h_1-h_0\|_\infty\le1$, the contraction estimate gives
\[
 \|h-h_n\|_\infty\le
 \sum_{j=n}^{\infty}\|h_{j+1}-h_j\|_\infty
 \le\sum_{j=n}^{\infty}q^j=E_n.
\]
Multiplying the pointwise errors by $r_-$ and $r_+$ proves
\eqref{eq:restored-support-gap-error}.
\end{proof}

\begin{remark}[Controlling unsampled directions]
\label{rem:restored-direction-grid}
Finite directional sampling also requires a bound between samples.
The function $h_n$ is the support function of the convex hull of
$\{f_u(0):|u|=n\}$, which lies in the disk of radius
$R_n=(1-q^n)/(1-q)$. Therefore $g_n$ is Lipschitz in chordal distance
on $S^1$, with constant at most
$L_n=(r_-+r_+)R_n+2$.
If every direction is within chordal distance $\eta$ of a sampled
direction, then $\min_j g_n(u_j)\ge L_n\eta$ proves convex contact.
The inequalities must use rigorous numerical enclosures when
implemented. Convex contact alone does not certify contact of the
fractal pieces; it is an outer approximation to connectedness.
\end{remark}

\begin{theorem}[Rigidity on the outer ceiling]
\label{thm:restored-outer-rigidity}
Suppose $r_-+r_+=1$. Then $K$ is connected if and only if both normalized
semilinear multipliers are real. In this case $K$ is a real interval and
$f_-(K)\cap f_+(K)$ is a singleton. Equivalently, the part of the
connectedness locus on $\rho=2$ consists precisely of these real systems.
\end{theorem}

\begin{proof}
Assume first that $K$ is connected, and put $D=\operatorname{diam}K>0$.
The level-$n$ cylinder cover has maximum diameter at most $q^nD$ and
total diameter
\[
 \sum_{|u|=n}\operatorname{diam}f_u(K)
   =D(r_-+r_+)^n=D.
\]
Hence $\mathcal H^1(K)\le D$.

We give the equality argument explicitly. Choose a diametral pair
$x,y\in K$. If $z\in K$ is not on the segment $[x,y]$, the three
strictly positive numbers
\[
 \begin{aligned}
 R_x&=\frac{|x-y|+|x-z|-|y-z|}{2},\\
 R_y&=\frac{|x-y|+|y-z|-|x-z|}{2},\\
 R_z&=\frac{|x-z|+|y-z|-|x-y|}{2}
 \end{aligned}
\]
are radii of pairwise disjoint open balls centered at $x,y,z$.
Connectedness implies that the distance from each center takes every
value between zero and the corresponding radius on the part of $K$
inside that ball. Since a distance function is $1$-Lipschitz,
\[
 \mathcal H^1(K)\ge R_x+R_y+R_z
   =\frac{|x-y|+|x-z|+|y-z|}{2}>D,
\]
a contradiction. Thus $K\subset[x,y]$, and connectedness gives
$K=[x,y]$.

It remains to identify the line. The average map
$G=(f_-+f_+)/2$ is a strict real-linear contraction fixing zero.
Convexity of the segment gives $G(K)\subset K$. Therefore
$0\in K$, and then $-1=f_-(0)$ and $1=f_+(0)$ belong to $K$.
Its supporting line is the real line. Each nonzero semilinear map
sends a nondegenerate real interval to a real interval, so its complex
multiplier is real.

Conversely, for real multipliers the real-slice criterion gives a
connected interval when $r_-+r_+=1$. The two subintervals cover $K$,
their lengths sum to the length of $K$, and they are proper
subintervals. Their intersection is consequently one point.
\end{proof}

\begin{proposition}[Area of the first-piece contact]
\label{prop:restored-contact-area}
Let $Q=f_-(K)\cap f_+(K)$ and let $m_2$ denote planar Lebesgue measure.
For every parameter,
\begin{equation}\label{eq:restored-contact-area}
 m_2(Q)=(r_-^2+r_+^2-1)m_2(K).
\end{equation}
Every point on the boundary of the connectedness locus, relative to
the strict contraction domain, satisfies
\[
 r_-+r_+\ge1,\qquad r_-^2+r_+^2\le1,
 \qquad m_2(Q)=0.
\]
If $r_-^2+r_+^2<1$, one also has $m_2(K)=0$.
\end{proposition}

\begin{proof}
Inclusion--exclusion and similarity scaling give
\[
 m_2(K)=r_-^2m_2(K)+r_+^2m_2(K)-m_2(Q),
\]
which proves the identity. The connectedness locus is relatively
closed, so a boundary point belongs to it and satisfies the outer
bound. The open set $r_-^2+r_+^2>1$ lies inside the connectedness locus
by the universal inner bound. Hence boundary points satisfy the
opposite weak inequality. If its left side is strictly less than
one, nonnegativity in \eqref{eq:restored-contact-area} forces both
measures to vanish; at equality the contact measure vanishes directly.
\end{proof}

Zero-area contact can be a nondegenerate segment. The exact families
below exhibit this on the full parameter-space boundary.
